\chapter*{\center \Large Abstract}

Deep reinforcement learning agents that learn through imagined trajectories within a learned world model improve sample efficiency over purely model-free methods. However, these imagination-based approaches are sensitive to reward scale: when value predictions are trained with standard mean squared error losses, critic predictions can diverge when imagined returns span several orders of magnitude. This report investigates whether the symlog transformation, a symmetric logarithmic function recently introduced in DreamerV3, can stabilize this training.

We implement a simplified Dreamer-like architecture with a learned transition model and an actor-critic policy trained on imagined rollouts. We compare three configurations: a model-free actor-critic baseline, an imagination-based actor-critic with standard MSE losses, and an imagination-based actor-critic with symlog-transformed losses. We evaluate these configurations on CartPole-v1 and LunarLander-v3 using five random seeds, yielding thirty training runs.

Without the symlog transformation, imagination-based training suffers from critic loss divergence (reaching values near $10^{15}$ in LunarLander) and performance collapse. The symlog transformation prevents this divergence, maintaining stable critic losses and performance. All configurations converge on the simpler CartPole-v1 environment, which confirms that the instability depends on reward scale.

These findings indicate that symlog-transformed losses are important for stable imagination-based training, highlighting the role of loss design in reinforcement learning.

\paragraph{Keywords:} deep reinforcement learning, symlog, world model, imagination-based training, actor-critic, loss function adaptation
